\section{Introduction}
\label{sec:introduction}

% MANUSCRIPT-NODE: INTRO-P01
Off-policy policy optimization improves a current policy by reusing behavior collected by earlier, stale, or heterogeneous policies. This paradigm is central to offline reinforcement learning, replay-based control, recommender systems, and model post-training because it reduces the need for fresh interaction and exploits experience accumulated across policy iterations~\cite{levine2020offline,fujimoto2019off}. Historical data contain two complementary signals. Positive updates increase the likelihood of observed successful actions, whereas negative feedback suppresses known failures and competing modes~\cite{peng2019advantage,wang2020critic}. The Positive-only endpoint, which removes every negative update, therefore yields a useful stability reference but can restrict the learner to attraction toward the empirical positive target. Balanced repulsion can instead reshape the decision boundary and move the aggregate equilibrium beyond that target. The central question is not whether negative feedback should be kept or discarded wholesale, but how to preserve its policy-improvement value without allowing repeated reuse to turn it into excessive repulsion.
% END-MANUSCRIPT-NODE: INTRO-P01

% MANUSCRIPT-NODE: INTRO-P02
The difficulty appears when negative actions are historical rather than freshly sampled. Offline logs, replay buffers, stale actors, and asynchronous trajectories can keep presenting the same action after the current learner has changed~\cite{schulman2017proximal,haarnoja2018soft}. A negative sample may initially lie near the policy and provide locally informative boundary information. Its update then pushes the learner away, increasing the action's learner-relative distance or surprisal while the stored negative label remains active. In Gaussian policies, the mean score can grow with standardized distance, so subsequent reuse produces progressively larger repulsion. In categorical policies, the direct-logit score is bounded, yet repeated negative updates continue to lower the selected action's log-odds and can drive its probability toward the support boundary. Historical reuse therefore creates a policy-family-dependent transition from useful local feedback to destructive far-field influence, rather than merely exposing the learner to one unusually large update.
% END-MANUSCRIPT-NODE: INTRO-P02

% MANUSCRIPT-NODE: INTRO-P03
Existing methods already regulate negative or off-policy updates in several principled ways. Positive-only objectives remove negative terms; global coefficients and clipping reduce their scale; trust regions and behavior constraints restrict policy movement; rarity-aware rules target low-probability events; and quality filtering removes selected data~\cite{schulman2017proximal,kumar2020conservative,kostrikov2021offline,peng2019advantage}. These controls can improve stability, but they do not by themselves explain why the same negative action can be useful near the learner and destructive in the far field after repeated reuse. The issue is difficult to identify in ordinary logs because low reward, large negative advantage, rarity, and learner-relative distance are typically correlated. We therefore construct matched controls that use matching context and semantic role, holding negative-advantage magnitude, sample count, base coefficient, and policy stage fixed to isolate policy remoteness. The source-identification protocol tests whether remoteness independently changes gradient scale; a separate intervention protocol tests whether the resulting far-field influence transmits into drift, support-boundary events, and task-performance collapse. This separation motivates an aggregate theory rather than another static weighting heuristic.
% END-MANUSCRIPT-NODE: INTRO-P03

% MANUSCRIPT-NODE: INTRO-P04
Repulsive Dynamics separates three layers that are often conflated. First, learner-relative remoteness determines the score response of a fixed historical action. Second, repeated reuse creates an asymmetry: positive updates make a successful action more compatible with the learner and attenuate its score, whereas negative updates make a rejected action more remote and preserve or amplify its response. Third, positive and negative contributions aggregate into a signed policy field. When positive mass dominates and the signed target remains inside the feasible mean-parameter space, the policy has a finite stable equilibrium beyond the Positive-only target. At the critical balance the field can produce persistent drift, while negative dominance makes any finite stationary point unstable. The terminal manifestation depends on the policy family: Gaussian mean runaway yields unbounded far-field scores, whereas categorical policies approach the simplex boundary with bounded per-sample logit scores. Negative feedback is therefore neither uniformly beneficial nor uniformly harmful; its value depends on aggregate balance and learner-relative location.
% END-MANUSCRIPT-NODE: INTRO-P04

% MANUSCRIPT-NODE: INTRO-P05
DRPO acts on the mechanism identified by the theory rather than treating every negative sample as toxic. Starting from the signed empirical actor field, it reweights the aggregate negative term by learner-relative distance or surprisal. Nearby negative actions retain substantial weight and can continue to shape local boundaries or suppress competing bad modes. As an action becomes increasingly remote, the weight decays nonlinearly, preventing the far-field tail from dominating the positive attraction. The exponential envelope is chosen for a precise tail property: under finite-order growth of the unweighted score-times-advantage contribution, the weighted contribution vanishes as remoteness increases. This guarantee does not assume that a sample's utility decays exponentially and does not establish a universal ranking over all tapers. Quality-based selection and learner-relative remoteness control are distinct axes; the experiments evaluate them and global scaling under matched conditions rather than presuming that one control must always win.
% END-MANUSCRIPT-NODE: INTRO-P05

% MANUSCRIPT-NODE: INTRO-P06
The empirical program follows the same decomposition as the theory. First, we ask whether far-field or rare negative updates become disproportionately influential in realistic policy learning. Second, using matched quality--distance and quality--rarity controls, we test whether remoteness independently changes negative influence and whether targeted far-field interventions interrupt drift or boundary events. Third, we test the predicted sequence from the Positive-only platform through stable extrapolation to persistent drift or instability, and compare selective attenuation with global and budget-matched controls. Fourth, we evaluate whether the resulting control improves external tasks. C-U1 and D-U1 provide controlled continuous and categorical identification; Hopper/D4RL and Countdown provide external validity and do not replace those controlled mechanisms. Every dynamics claim is terminal-audited, and task-performance collapse, support or variance-boundary events, and NaN/Inf numerical failure are reported separately. We contribute a layered theory of repeated repulsion, matched causal identification of learner-relative remoteness, DRPO control of the far-field negative term, and an evidence chain that keeps mechanism claims distinct from external task claims.
% END-MANUSCRIPT-NODE: INTRO-P06
